% Studio di Fattibilità
% da compilare con il comando pdflatex Studio_di_Fattibilita_x.x.x.tex

% Dichiarazioni di ambiente e inclusione di pacchetti
% da usare tramite il comando % Dichiarazioni di ambiente e inclusione di pacchetti
% da usare tramite il comando % Dichiarazioni di ambiente e inclusione di pacchetti
% da usare tramite il comando \input{../../util/hx-ambiente}

\documentclass[a4paper,titlepage]{article}
\usepackage[T1]{fontenc}
\usepackage[utf8]{inputenc}
\usepackage[english,italian]{babel}
\usepackage{microtype}
\usepackage{lmodern}
\usepackage{underscore}
\usepackage{graphicx}
\usepackage{eurosym}
\usepackage{float}
\usepackage{fancyhdr}
\usepackage[table]{xcolor}
\usepackage{longtable}
\usepackage{chngpage}
\usepackage{grffile}
\usepackage[titles]{tocloft}
\usepackage{hyperref}
\hypersetup{hidelinks}

\usepackage{../../util/hx-vers}
\usepackage{../../util/hx-macro}
\usepackage{../../util/hx-front}

% solo se si vuole una nuova pagina ad ogni \section:
\usepackage{titlesec}
\newcommand{\sectionbreak}{\clearpage}

% stile di pagina:
\pagestyle{fancy}

% solo se si vuole eliminare l'indentazione ad ogni paragrafo:
\setlength{\parindent}{0pt}

% intestazione:
\lhead{\Large{\proj}}
\rhead{\includegraphics[keepaspectratio=true,width=50px]{../../util/hivex_logo2.png}}
\renewcommand{\headrulewidth}{0.4pt}

% pie' di pagina:
\lfoot{\email}
\rfoot{\thepage}
\cfoot{}
\renewcommand{\footrulewidth}{0.4pt}

% spazio verticale tra le celle di una tabella:
\renewcommand{\arraystretch}{1.5}

% profondità di indicizzazione:
\setcounter{tocdepth}{4}
\setcounter{secnumdepth}{4}


\documentclass[a4paper,titlepage]{article}
\usepackage[T1]{fontenc}
\usepackage[utf8]{inputenc}
\usepackage[english,italian]{babel}
\usepackage{microtype}
\usepackage{lmodern}
\usepackage{underscore}
\usepackage{graphicx}
\usepackage{eurosym}
\usepackage{float}
\usepackage{fancyhdr}
\usepackage[table]{xcolor}
\usepackage{longtable}
\usepackage{chngpage}
\usepackage{grffile}
\usepackage[titles]{tocloft}
\usepackage{hyperref}
\hypersetup{hidelinks}

\usepackage{../../util/hx-vers}
\usepackage{../../util/hx-macro}
\usepackage{../../util/hx-front}

% solo se si vuole una nuova pagina ad ogni \section:
\usepackage{titlesec}
\newcommand{\sectionbreak}{\clearpage}

% stile di pagina:
\pagestyle{fancy}

% solo se si vuole eliminare l'indentazione ad ogni paragrafo:
\setlength{\parindent}{0pt}

% intestazione:
\lhead{\Large{\proj}}
\rhead{\includegraphics[keepaspectratio=true,width=50px]{../../util/hivex_logo2.png}}
\renewcommand{\headrulewidth}{0.4pt}

% pie' di pagina:
\lfoot{\email}
\rfoot{\thepage}
\cfoot{}
\renewcommand{\footrulewidth}{0.4pt}

% spazio verticale tra le celle di una tabella:
\renewcommand{\arraystretch}{1.5}

% profondità di indicizzazione:
\setcounter{tocdepth}{4}
\setcounter{secnumdepth}{4}


\documentclass[a4paper,titlepage]{article}
\usepackage[T1]{fontenc}
\usepackage[utf8]{inputenc}
\usepackage[english,italian]{babel}
\usepackage{microtype}
\usepackage{lmodern}
\usepackage{underscore}
\usepackage{graphicx}
\usepackage{eurosym}
\usepackage{float}
\usepackage{fancyhdr}
\usepackage[table]{xcolor}
\usepackage{longtable}
\usepackage{chngpage}
\usepackage{grffile}
\usepackage[titles]{tocloft}
\usepackage{hyperref}
\hypersetup{hidelinks}

\usepackage{../../util/hx-vers}
\usepackage{../../util/hx-macro}
\usepackage{../../util/hx-front}

% solo se si vuole una nuova pagina ad ogni \section:
\usepackage{titlesec}
\newcommand{\sectionbreak}{\clearpage}

% stile di pagina:
\pagestyle{fancy}

% solo se si vuole eliminare l'indentazione ad ogni paragrafo:
\setlength{\parindent}{0pt}

% intestazione:
\lhead{\Large{\proj}}
\rhead{\includegraphics[keepaspectratio=true,width=50px]{../../util/hivex_logo2.png}}
\renewcommand{\headrulewidth}{0.4pt}

% pie' di pagina:
\lfoot{\email}
\rfoot{\thepage}
\cfoot{}
\renewcommand{\footrulewidth}{0.4pt}

% spazio verticale tra le celle di una tabella:
\renewcommand{\arraystretch}{1.5}

% profondità di indicizzazione:
\setcounter{tocdepth}{4}
\setcounter{secnumdepth}{4}


\version{1.0.0}
\creaz{20 dicembre 2016}
\author{\AZ}
\supervisor{\LS}
\uso{interno}
\dest{\ALL}
\title{Studio di Fattibilità}
% \date{9 gennaio 2017}

\begin{document}
\maketitle
% diario delle modifiche per il piano di progetto
% da includere con % diario delle modifiche per il piano di progetto
% da includere con % diario delle modifiche per il piano di progetto
% da includere con \include{diario}

\begin{diario}
	1.0.1 & {\GG} (Responsabile) & 03/02/2017 & Correzione degli errori individuati nelle prime tre sezioni del documento \\ \hline
	1.0.0 & {\PB} (Responsabile) & 09/01/2017 & Approvazione documento \\ \hline
	0.1.0 & {\MM} (Verificatore) & 07/01/2017 & Verifica documento \\ \hline
	0.0.9 & {\PB} (Responsabile) & 05/01/2017 & Stesura Organigramma \\ \hline
	0.0.8 & {\LB} (Responsabile) & 05/01/2017 & Stesura Consuntivo di Periodo \\ \hline
	0.0.7 & {\LB} (Responsabile) & 04/01/2017 & Stesura Preventivo di Periodo \\ \hline
	0.0.6 & {\LB} (Responsabile) & 03/01/2017 & Inserimento Gantt Diagrammi delle Attività \\ \hline
	0.0.5 & {\PB} (Responsabile) & 29/12/2016 & Stesura Pianificazione delle Attività \\ \hline
	0.0.4 & {\PB} (Responsabile) & 28/12/2016 & Stesura Introduzione \\ \hline
	0.0.3 & {\LB} (Responsabile) & 27/12/2016 & Stesura Modello di sviluppo \\ \hline
	0.0.2 & {\PB} (Responsabile) & 27/12/2016 & Stesura Analisi dei Rischi \\ \hline
	0.0.1 & {\LB} (Responsabile) & 26/12/2016 & Stesura scheletro \\ \hline
\end{diario}

\begin{diario}
	1.0.1 & {\GG} (Responsabile) & 03/02/2017 & Correzione degli errori individuati nelle prime tre sezioni del documento \\ \hline
	1.0.0 & {\PB} (Responsabile) & 09/01/2017 & Approvazione documento \\ \hline
	0.1.0 & {\MM} (Verificatore) & 07/01/2017 & Verifica documento \\ \hline
	0.0.9 & {\PB} (Responsabile) & 05/01/2017 & Stesura Organigramma \\ \hline
	0.0.8 & {\LB} (Responsabile) & 05/01/2017 & Stesura Consuntivo di Periodo \\ \hline
	0.0.7 & {\LB} (Responsabile) & 04/01/2017 & Stesura Preventivo di Periodo \\ \hline
	0.0.6 & {\LB} (Responsabile) & 03/01/2017 & Inserimento Gantt Diagrammi delle Attività \\ \hline
	0.0.5 & {\PB} (Responsabile) & 29/12/2016 & Stesura Pianificazione delle Attività \\ \hline
	0.0.4 & {\PB} (Responsabile) & 28/12/2016 & Stesura Introduzione \\ \hline
	0.0.3 & {\LB} (Responsabile) & 27/12/2016 & Stesura Modello di sviluppo \\ \hline
	0.0.2 & {\PB} (Responsabile) & 27/12/2016 & Stesura Analisi dei Rischi \\ \hline
	0.0.1 & {\LB} (Responsabile) & 26/12/2016 & Stesura scheletro \\ \hline
\end{diario}

\begin{diario}
	1.0.1 & {\GG} (Responsabile) & 03/02/2017 & Correzione degli errori individuati nelle prime tre sezioni del documento \\ \hline
	1.0.0 & {\PB} (Responsabile) & 09/01/2017 & Approvazione documento \\ \hline
	0.1.0 & {\MM} (Verificatore) & 07/01/2017 & Verifica documento \\ \hline
	0.0.9 & {\PB} (Responsabile) & 05/01/2017 & Stesura Organigramma \\ \hline
	0.0.8 & {\LB} (Responsabile) & 05/01/2017 & Stesura Consuntivo di Periodo \\ \hline
	0.0.7 & {\LB} (Responsabile) & 04/01/2017 & Stesura Preventivo di Periodo \\ \hline
	0.0.6 & {\LB} (Responsabile) & 03/01/2017 & Inserimento Gantt Diagrammi delle Attività \\ \hline
	0.0.5 & {\PB} (Responsabile) & 29/12/2016 & Stesura Pianificazione delle Attività \\ \hline
	0.0.4 & {\PB} (Responsabile) & 28/12/2016 & Stesura Introduzione \\ \hline
	0.0.3 & {\LB} (Responsabile) & 27/12/2016 & Stesura Modello di sviluppo \\ \hline
	0.0.2 & {\PB} (Responsabile) & 27/12/2016 & Stesura Analisi dei Rischi \\ \hline
	0.0.1 & {\LB} (Responsabile) & 26/12/2016 & Stesura scheletro \\ \hline
\end{diario}
\tableofcontents



\section{Introduzione}
	\subsection{Scopo del documento}
	Lo scopo del documento è quello di presentare le motivazioni che hanno spinto il gruppo \hx{} a scegliere il capitolato d'appalto C6. 
	Inoltre sarà presente una descrizione degli altri capitolati con le ragioni che hanno portato il gruppo ad escluderli dalla scelta.
	
	\subsection{Scopo del prodotto}
	\scopo
	
	\subsection{Glossario}
	\presgloss
	
	\subsection{Riferimenti}
		\subsubsection{Normativi}
		\begin{itemize}
			\item \textbf{Norme di Progetto}: \NdP;
		\end{itemize}
		\subsubsection{Informativi}
		\begin{itemize}
			\item \textbf{Glossario}: \Glossario;
			\item \textbf{Capitolato d'Appalto C1}, \emph{APIM}: An \gloss{API} Market Platform \\
			\url{www.math.unipd.it/~tullio/IS-1/2016/Progetto/C1.pdf}, visitato il 08/02/2017;
			\item \textbf{Capitolato d'Appalto C2}, \emph{AtAVi}: Accoglienza tramite Assistente Virtuale \\
			\url{www.math.unipd.it/~tullio/IS-1/2016/Progetto/C2.pdf}, visitato il 08/02/2017;
			\item \textbf{Capitolato d'Appalto C3}, \emph{DeGeOP}: A Designer and Geo-localizer Web App for Organizational Plants \\
			\url{www.math.unipd.it/~tullio/IS-1/2016/Progetto/C3.pdf}, visitato il 08/02/2017;
			\item \textbf{Capitolato d'Appalto C4}, \emph{eBread}: applicazione di lettura per dislessici \\
			\url{www.math.unipd.it/~tullio/IS-1/2016/Progetto/C4.pdf}, visitato il 08/02/2017;
			\item \textbf{Capitolato d'Appalto C5}, \emph{Monolith}: an interactive bubble provider \\
			\url{www.math.unipd.it/~tullio/IS-1/2016/Progetto/C5.pdf}, visitato il 08/02/2017;
			\item \textbf{Capitolato d'Appalto C6}, \emph{SWEDesigner}: editor di diagrammi UML con generazione di codice \\
			\url{www.math.unipd.it/~tullio/IS-1/2016/Progetto/C6.pdf}, visitato il 08/02/2017;
		\end{itemize}


	
\section{Studio di fattibilità del capitolato scelto}
	\subsection{\proj: editor di diagrammi UML con generazione di codice}
		\subsubsection{Descrizione}
		Lo scopo principale del capitolato è quello di realizzare un'\gloss{applicazione web} composta da almeno due disegnatori
		ed un coordinatore, che permetta all'utente di costruire diagrammi UML rappresentativi del prodotto da svilupparsi e generi a partire da essi codice Java o JavaScript.
		
		In particolare l'editor dovrebbe concentrarsi su uno specifico dominio applicativo (l'esempio fornito nel capitolato d'appalto è relativo ai software di giochi da tavolo) per andare a costituire una sorta di \gloss{framework} per la produzione della relativa famiglia di prodotti.

		Lo scopo profondo del capitolato è quello di analizzare quali diagrammi UML siano i più adatti allo scopo di descrivere l'architettura 
		e le componenti funzionali di un progetto software e come si possa concertare una loro integrazione per arrivare a produrre un eseguibile vero e proprio. 

		Attualmente, infatti, UML non è altro che un grande aggregato di convenzioni di modellazione che coprono, sì, a diversi livelli di raffinamento differenti fasi e aspetti dello sviluppo di prodotti software, ma al contempo risultano a tal punto eterogenei l'uno rispetto all'altro da non risultare integrabili e componibili. 

		La richiesta è quindi quella di indagare come si potrebbe superare la debolezza intrinseca dell'Unified Modeling Language nella sincronizzazione tra codice e quanto disegnato nei diagrammi stessi e sperimentare la possibilità di avvicinare le due fasi di progettazione e codifica, esplorando le interazioni e le eventuali estensioni necessarie a legare fra loro queste due fasi.
		\subsubsection{Studio del dominio}
			\paragraph{Dominio applicativo} Il dominio applicativo del seguente capitolato riguarda l'ambito della generazione di codice a partire dal linguaggio UML. Il prodotto si rivolge quindi, ad un target di programmatori interessati alla creazione di programmi appartenenti all'ambito dei giochi da tavolo. I clienti verranno guidati durante il percorso di realizzazione del proprio prodotto, ma dovranno possedere una certa esperienza e conoscenza nell'ambito della programmazione software, del linguaggio UML e dei suoi diagrammi.
			
			\paragraph{Dominio tecnologico}
			Il capitolato richiede che il sistema sia realizzato con tecnologie Web. In particolare la parte \gloss{server} deve essere realizzata in Java con il server 
			\gloss{Tomcat} o in Javascript con il server \gloss{Node.js}. La parte \gloss{client}, invece, deve essere eseguibile in un browser HTML5 ed utilizzare fogli di stile CSS3 per l'aspetto estetico e JavaScript per la parte attiva.
		
		\subsubsection{Aspetti positivi e negativi}
		Molti sono stati riconosciuti dal gruppo come gli aspetti positivi di tale progetto:
		\begin{itemize}
			\item le problematiche affrontate dal capitolato hanno suscitato in tutti i membri una forte motivazione, legata alla necessità di confrontarsi 
			con problemi quali la generazione automatica di codice e il raffinamento dell'UML;
			\item la chiarezza e precisione con cui le caratteristiche e i requisiti del progetto sono stati elencati e descritti;
			\item la possibilità di collaborare attivamente con quella che è considerata 
			la prima \emph{software house} italiana;
			\item la possibilità di operare con un insieme di tecnologie assodato, parte integrante del programma di studio e di ampia diffusione.
		\end{itemize}
		Fra le criticità discusse spiccano per importanza:
		\begin{itemize}
			\item la vastità e la complessità del problema da affrontare, certamente non banale e necessitante di una riflessione 
			estremamente approfondita e ponderata;
			\item la necessità di conoscere al meglio i costrutti e le astrazioni del linguaggio di generazione del codice, 
			in modo tale da garantire una buona integrazione fra le due fasi di architettura e scrittura di codice.
		\end{itemize}
		\subsubsection{Valutazione finale} 
		Il fatto di creare strumenti che regimentino il processo di sviluppo e permettere di fattorizzare gli aspetti comuni ad una certa famiglia di 
		prodotti software, unito all'esplorazione della generazione di codice in uno specifico linguaggio di programmazione, a partire dalla rappresentazione 
		dell'architettura software, hanno suscitato un elevato grado di interesse da parte dei membri del gruppo. Quest'ultimo, combinato con un insieme non 
		trascurabile di aspetti positivi, ha fatto sì che il gruppo \hx{} scegliesse tale capitolato d'appalto.
	
\newpage

\section{Studio di fattibilità degli altri capitolati}
	\subsection{Capitolato C1: \emph{APIM}: An API Market Platform}
		\subsubsection{Descrizione}
		Lo scopo di questo capitolato è quello di creare un'applicazione web che rappresenti una sorta di \gloss{marketplace} per la consultazione e 
		condivisione di microservizi. Tra le funzionalità che il prodotto software dovrebbe fornire vi sono la possibilità di registrare le API 
		di un microservizio, consultarle e ritrovare la relativa documentazione e visualizzare i relativi dati tecnici; inoltre è richiesta la 
		funzionalità di associare ad ogni API diverse chiavi d'uso, monitorare il suo utilizzo, limitare il suo uso da parte di utenti in possesso 
		di chiavi scadute o non valide.
		\subsubsection{Studio del dominio}
			\paragraph{Dominio applicativo}
			Come stanno facendo anche grandi aziende operanti nel settore della tecnologia digitale, il team di ItalianaSoftware sta focalizzando i propri sforzi su una nuova visione dell'architettura dei programmi basata non su grandi moduli funzionali bensì su unità minimali: i microservizi. Questi sono in grado di comporsi e cooperare in modo tale da formare aggregati sempre più complessi fino ad arrivare ad un'unica applicazione monolitica. Un'innovazione, questa, che potrebbe trasformare completamente la struttura dei sistemi informativi: da monolitici, con processi e dati difficilmente modificabili, diventerebbero adattivi e flessibili in quanto composti da tanti componenti interdipendenti, ciascuno dei quali offre funzionalità ben delineate.
			\paragraph{Dominio tecnologico}
			Per quanto riguarda il set di tecnologie da utilizzarsi, ItalianaSoftware richiede ai gruppi impegnati sul progetto:
			\begin{itemize}
			\item di sviluppare il front end della web application con l'ausilio della terna JavaScript, HTML e CSS3;
			\item di scegliere tra una base di dati \gloss{SQL} o \gloss{NoSQL} per quanto riguarda il database.
			\end{itemize}
			L'azienda suggerisce fortemente di utilizzare per la rappresentazione delle interfacce e per la creazione del marketplace
			il linguaggio orientato ai microservizi: Jolie.
		\subsubsection{Aspetti positivi e negativi}
		Fra le note positive emerse nell'analisi del capitolato ci sono:
		\begin{itemize}
			\item la possibilità di confrontarsi con quella che potrebbe essere a tutti gli effetti una vera e propria rivoluzione nell'intendere 
			l'architettura dei sistemi informativi, composti da moduli unitari componibili, riusabili e scalabili;
			\item la disponibilità dell'azienda a collaborare attivamente con i team impegnati nel progetto;
			\item l'opportunità di apprendere il linguaggio Jolie direttamente dall'azienda che lo ha creato.
		\end{itemize}
		Uno dei fattori chiave nella decisione se partecipare o meno all'appalto di tale progetto è proprio l'uso del linguaggio Jolie.

		Da un lato la necessità di imparare un nuovo linguaggio di programmazione ha sempre un importante valore a livello formativo, dall'altro, però, non è ancora ben chiaro quali siano le reali potenzialità di Jolie per affermarsi nel sovraffollato panorama tecnologico odierno.
		\subsubsection{Valutazione finale}
		Il gruppo \hx{} ha ritenuto opportuno scegliere un capitolato che permettesse il confronto con tecnologie che per i membri fossero nuove, ma al contempo già pienamente affermate in quanto a utilizzo, funzionalità, potenzialità. Tale scelta non trova riscontro in questo capitolato.
		
	\subsection{Capitolato C2: \emph{AtAVi}: Accoglienza tramite Assistente Virtuale}
		\subsubsection{Descrizione}
		Lo scopo di questo capitolato è quello di realizzare un applicativo web che permetta una prima accoglienza degli ospiti in visita all'ufficio dell'azienda zero12. Il progetto dovrebbe comporsi di tre parti ben distinte: 
		\begin{itemize}
			\item un'interfaccia web per permettere all'utente di interagire con l'utilizzatore del sistema;
			\item un'interfaccia per trasmettere informazioni attraverso il canale comunicativo aziendale \gloss{Slack};
			\item un sistema che sfrutti gli \gloss{AWS} per l'interazione con le API dell'assistente virtuale scelto.
		\end{itemize}
		Il capitolato si configura pertanto sia come un'indagine su ulteriori ambiti di integrazione di tali tecnologie, sia come la dimostrazione di come anche in ambito aziendale esse possano contribuire a rendere più efficiente e moderno l'ambiente lavorativo.
		\subsubsection{Studio del dominio}
			\paragraph{Dominio applicativo}
			Il progetto proposto si inserisce all'interno di un quadro di ampio respiro quale la trasformazione digitale dei sistemi aziendali e la presenza sempre più pregnante nella vita quotidiana di applicativi basati proprio sull'interazione con assistenti virtuali, in grado di organizzare e gestire i nostri impegni e adattarsi alle nostre preferenze e richieste.
			
			\paragraph{Dominio tecnologico}
			Fra le tecnologie con cui è richiesto confrontarsi per la realizzazione di tale progetto vi sono:
			\begin{itemize}
				\item i servizi offerti dall'\gloss{infrastruttura} degli AWS;
				\item database NoSQL per la memorizzazione e l'analisi delle informazioni raccolte nel corso delle interazioni con il cliente;
				\item linguaggio di programmazione Node.js;
				\item il sistema di comunicazione aziendale Slack;
				\item gli \gloss{SDK} dei principali assistenti virtuali sul mercato, quali Cortana, Siri e Google Home.
			\end{itemize}
		\subsubsection{Aspetti positivi e negativi}
		Fra i lati positivi del capitolato figurano:
		\begin{itemize}
			\item l'offerta di formazione da parte dell'azienda sulle principali tecnologie;
			\item la possibilità di lavorare con il supporto di un server dedicato configurato tramite gli AWS.
		\end{itemize}
		Le criticità piu importati riguardano:
		\begin{itemize}
			\item lo scopo del capitolato, che consiste più nel trovare un modo per concertare e coordinare fra loro diversi sistemi 
			esistenti più che nell'elaborare qualcosa di effettivamente nuovo nell'ambito applicativo;
			\item la necessità di confrontarsi con un vasto assortimento di tecnologie non ben conosciute e di indubbia complessità;
			\item l'insufficiente chiarezza a livello dell'architettura software da realizzare, considerata particolarmente problematica da 
			parte di alcuni membri del gruppo.
		\end{itemize}
		\subsubsection{Valutazione finale}
		Il gruppo ritene che la mancanza di esperienza possa togliere ampio spazio ad una riflessione più approfondita sulle funzionalità e sulla 
		struttura del progetto stesso, per questo ha deciso di scartare tale capitolato d'appalto.
		
	\subsection{Capitolato C3: \emph{DeGeOP}: A Designer and Geo-localizer Web App for Organizational Plants}
		\subsubsection{Descrizione}
		Lo scopo del capitolato in esame è quello di creare un'interfaccia di tipo web app utilizzabile anche su mobile, per visualizzare su di una 
		mappa geografica gli elementi costitutivi del processo produttivo di una particolare azienda.

		Tale web app deve permettere di:
		\begin{itemize}
			\item definire la struttura organizzativa di un'azienda;
			\item riportare tale struttura a livello di mappa geografica;
			\item disegnare gli scenari di danno che possono colpire l'azienda;
			\item ottenere dei risultati di analisi relativamente a tali scenari di rischio.
		\end{itemize}
		Quest'ultimo punto verrebbe attuato utilizzando il database per le situazioni di danno dell'azienda e il server di analisi, sempre di proprietà dell'azienda,
		che opera utilizzando un algoritmo proprietario.

		L'applicazione deve inoltre garantire un corretto funzionamento su desktop, tablet e dispositivi mobile e utilizzare API e database proprietari dell'azienda.
		\subsubsection{Studio del dominio}
			\paragraph{Dominio applicativo}
			Il contesto in cui l'applicazione opera costituisce una realtà di attuale interesse, soprattutto in un paese come il nostro dove sempre più spesso le catastrofi climatiche, quali terremoti, inondazioni, trombe d'aria, mettono in difficoltà intere aziende.

			Data pertanto la mancanza di una piena consapevolezza riguardo a tale problematica (il rischio relativo alle interruzioni di lavoro è poco conosciuto) e la frequenza con cui questa si presenta, il contributo di tale applicazione potrebbe essere in futuro estremamente significativo.

			Il servizio offerto dall'applicativo andrebbe inoltre a coprire un vuoto: le aziende di assicurazioni, infatti, non hanno ora strumenti adeguati per valutare il rischio legato all'interruzione dei lavori per cause esterne. Molte aziende infatti si appoggiano al servizio RiskApp, alcune delle quali di notevole importanza nel ramo assicurazioni.
			\paragraph{Dominio tecnologico}
			Nonostante la relativa libertà nell'utilizzo dello \gloss{stack} tecnologico le tecnologie proposte sono:
			\begin{itemize}
				\item gli AWS come servizio per l'organizzazione di un server entro cui far operare l'applicativo;
				\item \gloss{Django}, come framework in \gloss{Python} per la realizzazione del front-end dell'applicativo;
				\item \gloss{PostgreSQL}, per la realizzazione del database di appoggio del software;
				\item Librerie e framework quali \gloss{Bootstrap} e \gloss{React.js} per la creazione di interfacce utente, hammer.js e Yeoman;
				\item Html, Css e JavaScript per la realizzazione dell'interfaccia utente.
			\end{itemize}
		\subsubsection{Aspetti positivi e negativi}
		Per quanto riguarda gli aspetti positivi, oltre all'utilità dell'applicazione stessa, non ne sono stati evidenziati altri.
		
		Tra le criticità trovate, spiccano:
		\begin{itemize}
			\item la poca chiarezza nella descrizione del progetto;
			\item la nebulosità dei requisiti del capitolato.
		\end{itemize}
		\subsubsection{Valutazione finale}
		Nessun membro del gruppo ha dimostrato particolare interesse o motivazione per il progetto in esame. Pertanto si è deciso di non procedere 
		troppo oltre nell'esame delle caratteristiche di tale proposta, così da consentire un'analisi più dettagliata di altri capitolati in grado 
		di stimolare una maggiore motivazione.
		
	\subsection{Capitolato C4: \emph{eBread}: applicazione di lettura per dislessici}
		\subsubsection{Descrizione}
		Il progetto prevede la realizzazione di un'applicazione per dispositivi mobili in ambiente Android che agevoli, grazie all'uso di appropriate tecnologie, la lettura da parte di persone affette da dislessia.

		La proposta richiede inoltre che l'applicativo, indifferentemente un lettore di ebook o un client di messaggistica, utilizzi un motore di sintesi vocale come ausilio alla lettura, evidenziando le parole del testo, in modo sincronizzato con l'audio sintetizzato. 

		Un'ulteriore funzionalità da fornire, considerata molto utile per persone affette da dislessia, consiste nella la possibilità di modificare la visualizzazione del contenuto da leggere attraverso l'uso di particolari sistemi di colori, spaziature fra lettere, differenti dimensioni 
		e tipo dei caratteri e layout di visualizzazione.
		\subsubsection{Studio del dominio}
			\paragraph{Dominio applicativo}
			Le problematiche che il progetto proposto dal capitolato si prevede di affrontare sono sicuramente di grande attualità.

			Da un lato la dislessia, che nel nostro paese al giorno d'oggi rappresenta il disturbo dell'apprendimento più diffuso e che incide non poco nella vita delle persone che ne sono affette in quanto rende difficoltosa una delle attività fondamentali della routine quotidiana: la lettura. 

			Dall'altro si spinge il gruppo fornitore a confrontarsi con una delle tecnologie che sta avendo negli ultimi tempi una diffusione sempre maggiore, quella del \gloss{Text-to-Speech}, alla base di applicazioni di grande importanza quali screenreader, lettori di messaggi, voci di assistenti virtuali e probabilmente futura interfaccia per tutte quelle attività in cui è precluso, o comunque limitato, l'utilizzo della vista.
			\paragraph{Dominio tecnologico}
			Viene lasciato un ampio margine di libertà per quanto riguarda le tecnologie da impiegarsi, purché risultino adeguate allo scopo. 

			Gli unici due vincoli principali in questo senso sono:
			\begin{itemize}
			\item la realizzazione di una applicazione per dispositivi mobili, preferibilmente per sistema operativo Android;
			\item l'utilizzo di un servizio di sintesi vocale, con consiglio per il motore di sintesi “Flexible and Adaptive Text-To-Speech”.
			\end{itemize}
			Nel caso si proceda allo sviluppo di un ebook reader, il proponente consiglia di utilizzare una libreria di terze parti preferibilmente 
			open source per accedere al contenuto. Allo stesso modo, nel caso in cui si decida di estendere le funzionalità di un client di messaggistica, 
			si suggerisce di partire da uno già esistente, possibilmente open source.
		\subsubsection{Aspetti positivi e negativi}
		Tra le criticità è stato evidenziato il fatto che in tale capitolato ci sia da sviluppare un applicazione per dispositivi mobili.

		Le applicazioni per dispositivi mobili, infatti, richiedono un particolare bagaglio di conoscenze e una certa esperienza in tale campo che il gruppo ha ritenuto di non possedere e di non essere in grado di apprendere in quantità necessaria per svolgere tale capitolato che si è prefissato.
		
		\subsubsection{Valutazione finale}
		Complice una presentazione del progetto forse sottotono ed una tematica, anche se attuale e viva soprattutto negli ambienti educativi della società odierna, non capace di spingere la motivazione dei membri del gruppo, si è scelto di procedere all'appalto di un differente capitolato.
		
	\subsection{Capitolato C5: \emph{Monolith}: an interactive bubble provider}
		\subsubsection{Descrizione}
		Lo scopo del capitolato è quello di realizzare un framework che consenta agli sviluppatori di creare in modo relativamente semplice delle 
		interactive bubbles in grado di integrarsi con l'ambiente l'applicazione di messagistica Rocket.chat.

		Una interactive bubble viene intesa dal proponente come una sorta di mini-applicativo autonomo all'interno del canale di comunicazione 
		che consente di arricchire di funzionalità il client di messaggistica in uso e condividere informazioni complesse e strutturate. 

		A scopo dimostrativo occorre includere delle bolle interattive predefinite pronte da essere utilizzate da parte degli utenti finali.
		
		\subsubsection{Studio del dominio}
			\paragraph{Dominio applicativo}
			Le problematiche affrontate dal presente capitolato risultano sicuramente interessanti ed attuali.

			Si riconosce prima di tutto che al giorno d'oggi il più importante canale comunicativo è costituito dalle applicazioni di chat, utilizzate in ogni ambito della vita quotidiana: per relazionarci con gli amici, per collaborare in ambito lavorativo, per rimanere 
			connessi a distanza.

			Inoltre sempre più stringente è la necessità di scambiarsi all'interno di tali contesti informazioni più complesse e strutturate di singoli messaggi (come ad esempio dati relativi ad un volo aereo, alla prenotazione in un certo locale, etc.) che rimangano aggiornate e coerenti con le fonti da cui sono state tratte.

			Le interactive bubbles rappresenterebbero la soluzione ideale a tale problematica, che al giorno d'oggi si risolve invece con uno scambio intensivo di numerosi messaggi singoli. 
			\paragraph{Dominio tecnologico}
			Le tecnologie da utilizzarsi nel progetto sono le seguenti:
			\begin{itemize}
				\item Javascript versione 6 per la versione dimostrativa da mostrare al proponente e il framework finale;
				\item framework quali \gloss{Angular.js} e React.js sono consigliati per la realizzazione del front-end dell'applicazione;
				\item \gloss{SCSS} per l'interfaccia utente.
			\end{itemize}
			Viene inoltre richiesto che il codice sorgente di Monolith e la relativa versione dimostrativa siano pubblicate e versionate utilizzando 
			\gloss{GitHub} o \gloss{Bitbucket}.
		\subsubsection{Aspetti positivi e negativi}
		L'ambito applicativo del progetto risulta sicuramente ricco di spunti interessanti, soprattutto per l'importanza cruciale che i client di messaggistica odierni hanno assunto nello scambio di informazioni.

		Tra gli aspetti negativi messi in evidenza ci sono:
		\begin{itemize}
			\item lontananza dei proponenti, la quale potrebbe incidere su una comunicazione proficua con l'azienda proponente;
			\item lo sviluppo di un'applicazione per dispositivi mobili e tutte le criticità che ne derivano;
			\item un insieme di tecnologie ritenute non molto interessanti dai membri del gruppo.
		\end{itemize}
		\subsubsection{Valutazione finale}
		La maggior parte dei membri del gruppo non ha dimostrato grande interesse in merito alla proposta, e pertanto si è scelto di procedere all'appalto di un differente capitolato.

\end{document}
